\vkthmsync{13}
\paragraph{\trTriangulated\trACategory。}\label{sub:triangulated-categories}
称一个\trACategory{}是\textbf{\trTriangulated}（triangulated）的，若它非空%
（$\Ob\scr{A}\ne\emptyset$）且
\begin{itemize}
\item $H^0(\scr{A})$ 中的任何\trMorphism{} $[c_1]$ 都可以补全成一个%
\trExactTriangle{} \eqref{eq:exact-triangle-diagram}；
\item 对任何\trObject{} $Y$，都存在某个 $\widetilde{Y}$，使得在
$H^0(\scr{A})$ 中 $S\widetilde{Y}\cong Y$。
\end{itemize}
对此稍作解释是合适的。由\cref{thm:cone-shift-sum-relation}，第一条条件蕴含%
每个\trObject{} $Y$ 都有\trSuspension{} $SY$，这使第二条条件变得有意义。若换一种
更侧重\trFunctor{}的说法：第一条条件使我们能够构造一个\trShiftFunctor{}%
$S\colon\scr{A}\to\scr{A}$（见\cref{sub:shift}），而第二条条件则断言这是一个%
\trQuasiequivalence。

\begin{example}\label{thm:mod-categories-triangulated}
\trACategory{} $\scr{Q}=\mathit{mod}(\scr{A})$ 都是\trTriangulated{}的。
\end{example}

\begin{proposition}\label{thm:h0-triangulated-category}
若\trACategory{} $\scr{A}$ 是\trTriangulated{}的，则 $H^0(\scr{A})$ 是古典意义
下的\trTriangulated\trCategory。进而，若\trACategory{} $\scr{A}$、$\scr{B}$ 是%
\trTriangulated{}的，而 $\scr{F}\colon\scr{A}\to\scr{B}$ 是\trAFunctor，则
$H^0(\scr{F})$ 是\trExact{}\trFunctor。
\end{proposition}

我们回忆，（古典意义下的）\trTriangulated\trCategory{} [141], [61] 是一个%
\trAdditive\trCategory，连同“\trTranslation”（translation）自\trEquivalence{} $T$
与一族“\trDistinguishedTriangle”（distinguished triangles），后者是形如下式的%
\trPeriodic{}性图表
\begin{equation*}
\cdots\longrightarrow Y_0\longrightarrow Y_1\longrightarrow Y_2
\longrightarrow TY_0\longrightarrow TY_1\longrightarrow TY_2
\longrightarrow\cdots.
\end{equation*}
在我们的情形，\trTranslation{}由\trShiftFunctor{}诱导，即 $T=H^0(S)$，而%
\trDistinguishedTriangle{}就是\trExactTriangle，其中我们用
\eqref{eq:shift-hom-isomorphism} 把 $[c_3]$ 化成一个\trMorphism{}%
$Y_2\to SY_0$。
